\section{B}

% baden (to bathe)
\begin{verbentry}{
  style = {default},
  toc = {baden (to bathe)},
  name = {baden},
  english = {to bathe},
  type = {regular},
  auxiliary = {haben}
}


\vspace{5pt}
\hrule
\vspace{5pt}

\begin{tabular}{rl}
\multicolumn{2}{l}{\textbf{PRÄSENS}}\\
ich & \textbf{bade}\\
du & \textbf{badest}\\
er/sie/es & \textbf{badet}\\
wir & \textbf{baden}\\
ihr & \textbf{badet}\\
sie/Sie & \textbf{baden}\\
\end{tabular}
\hfill
\begin{tabular}{rl}
\multicolumn{2}{l}{\textbf{PRÄTERITUM}}\\
ich & \textbf{badete}\\
du & \textbf{badetest}\\
er/sie/es & \textbf{badete}\\
wir & \textbf{badeten}\\
ihr & \textbf{badetet}\\
sie/Sie & \textbf{badeten}\\
\end{tabular}
\hfill
\begin{tabular}{rl}
\multicolumn{2}{l}{\textbf{IMPERATIV}}\\
& \textbf{bade (du)!}\\
& \textbf{baden wir!}\\
& \textbf{badet (ihr)!}\\
& \textbf{baden Sie!}\\
\end{tabular}

\vspace{10pt}

\begin{tabular}{rl}
\multicolumn{2}{l}{\textbf{PERFEKT}}\\
& haben + \textbf{gebadet}\\
\end{tabular}
\hfill
\begin{tabular}{rl}
\multicolumn{2}{l}{\textbf{FUTURE I}}\\
& werden + \textbf{baden}\\
\end{tabular}
\hfill
\begin{tabular}{rl}
\multicolumn{2}{l}{\textbf{PLUSQUAMPERFEKT}}\\
& hatten + \textbf{gebadet}\\
\end{tabular}

\vspace{10pt}

\begin{tabular}{lll}
\multicolumn{3}{l}{\textbf{MIT PRÄPOSITIONEN}}\\
& \textbf{baden in} + Dat & (to bathe / swim in)\\
\end{tabular}
\hfill
\begin{tabular}{lll}
\multicolumn{3}{l}{\textbf{TRENNBARE VERBEN}}\\
& \textbf{---} & \\
\end{tabular}
\vspace{-5pt}
\end{verbentry}

% bauen (to build)
\begin{verbentry}{
  style = {default},
  toc = {bauen (to build)},
  name = {bauen},
  english = {to build},
  type = {regular},
  auxiliary = {haben}
}
 
    
     \vspace{5pt}

    \hrule
    
    \vspace{5pt}

  \begin{tabular}{rl}
     \multicolumn{2}{l}{\textbf{PRÄSENS} }\\
    ich & \textbf{baue}\\
    du & \textbf{baust}\\
    er/sie/es & \textbf{baut}\\
    wir & \textbf{bauen}\\
    ihr & \textbf{baut}\\
    sie/Sie & \textbf{bauen}\\
  \end{tabular}
  \hfill
     \begin{tabular}{rl}
    \multicolumn{2}{l}{\textbf{PRÄTERITUM} }\\
    ich & \textbf{baute}\\
    du & \textbf{bautest}\\
    er/sie/es & \textbf{baute}\\
    wir & \textbf{bauten}\\
    ihr & \textbf{bautet}\\
    sie/Sie & \textbf{bauten}\\
  \end{tabular}
  \hfill
  \begin{tabular}{rl}
     \multicolumn{2}{l}{\textbf{IMPERATIV} }\\
   & \textbf{bau(e) (du)!}\\
   & \textbf{bauen wir!}\\
   & \textbf{baut (ihr)!}\\
   & \textbf{bauen Sie!}\\
   & \vspace{10pt}
  \end{tabular}

    \vspace{10pt}

 \begin{tabular}{rl}
     \multicolumn{2}{l}{\textbf{PERFEKT}}\\
   & haben + \textbf{gebaut}\\
  \end{tabular}
\hfill
   \begin{tabular}{rl}
   
    \multicolumn{2}{l}{\textbf{FUTURE I} }\\
   & werden + \textbf{bauen}\\
 
  \end{tabular}
\hfill
   \begin{tabular}{rl}
      \multicolumn{2}{l}{\textbf{PLUSQUAMPERFEKT}}\\
   & hatten + \textbf{gebaut}\\
  \end{tabular}

   \vspace{10pt}

   \begin{tabular}{lll}
      \multicolumn{3}{l}{\textbf{MIT PRÄPOSITIONEN}}\\
& \textbf{bauen auf} + Akk & (to rely on / build on)\\
& \textbf{bauen an} + Dat & (to be working on building)\\
& \textbf{bauen in} + Akk & (to install into)\\
   \end{tabular}
  \hfill
   \begin{tabular}{lll}
      \multicolumn{3}{l}{\textbf{TRENNBARE VERBEN}}\\
 & \textbf{ab$|$bauen} & (to dismantle / reduce)\\
 & \textbf{auf$|$bauen} & (to build up / establish)\\
 & \textbf{ein$|$bauen} & (to install)\\
 & \textbf{um$|$bauen} & (to rebuild / remodel)
  \end{tabular}

\vspace{-5pt}
\end{verbentry}

% beantragen (to apply for)
\begin{verbentry}{
  style = {acc},
  toc = {beantragen (to apply for)},
  name = {beantragen},
  english = {to apply for},
  type = {regular, + Akkusativ},
  auxiliary = {haben}
}


\vspace{5pt}
\hrule
\vspace{5pt}

\begin{tabular}{rl}
\multicolumn{2}{l}{\textbf{PRÄSENS}}\\
ich & \textbf{beantrage}\\
du & \textbf{beantragst}\\
er/sie/es & \textbf{beantragt}\\
wir & \textbf{beantragen}\\
ihr & \textbf{beantragt}\\
sie/Sie & \textbf{beantragen}\\
\end{tabular}
\hfill
\begin{tabular}{rl}
\multicolumn{2}{l}{\textbf{PRÄTERITUM}}\\
ich & \textbf{beantragte}\\
du & \textbf{beantragtest}\\
er/sie/es & \textbf{beantragte}\\
wir & \textbf{beantragten}\\
ihr & \textbf{beantragtet}\\
sie/Sie & \textbf{beantragten}\\
\end{tabular}
\hfill
\begin{tabular}{rl}
\multicolumn{2}{l}{\textbf{IMPERATIV}}\\
& \textbf{beantrage (du)!}\\
& \textbf{beantragen wir!}\\
& \textbf{beantragt (ihr)!}\\
& \textbf{beantragen Sie!}\\
\end{tabular}

\vspace{10pt}

\begin{tabular}{rl}
\multicolumn{2}{l}{\textbf{PERFEKT}}\\
& haben + \textbf{beantragt}\\
\end{tabular}
\hfill
\begin{tabular}{rl}
\multicolumn{2}{l}{\textbf{FUTURE I}}\\
& werden + \textbf{beantragen}\\
\end{tabular}
\hfill
\begin{tabular}{rl}
\multicolumn{2}{l}{\textbf{PLUSQUAMPERFEKT}}\\
& hatten + \textbf{beantragt}\\
\end{tabular}

\vspace{10pt}

\begin{tabular}{lll}
\multicolumn{3}{l}{\textbf{MIT PRÄPOSITIONEN}}\\
& \textbf{beantragen bei} + Dat & (to apply at)\\
\end{tabular}
\hfill
\begin{tabular}{lll}
\multicolumn{3}{l}{\textbf{TRENNBARE VERBEN}}\\
& \textbf{---} & \\
\end{tabular}
\vspace{-5pt}
\end{verbentry}

% bedeuten (to mean something)
\begin{verbentry}{
  style = {acc},
  toc = {bedeuten (to mean something)},
  name = {bedeuten},
  english = {to mean something},
  type = {regular, + Akkusativ},
  auxiliary = {haben}
}
 
    
     \vspace{5pt}

    \hrule
    
    \vspace{5pt}

  \begin{tabular}{rl}
     \multicolumn{2}{l}{\textbf{PRÄSENS} }\\
    ich & \textbf{bedeute}\\
    du & \textbf{bedeutest}\\
    er/sie/es & \textbf{bedeutet}\\
    wir & \textbf{bedeuten}\\
    ihr & \textbf{bedeutet}\\
    sie/Sie & \textbf{bedeuten}\\
  \end{tabular}
  \hfill
     \begin{tabular}{rl}
    \multicolumn{2}{l}{\textbf{PRÄTERITUM} }\\
    ich & \textbf{bedeutete}\\
    du & \textbf{bedeutetest}\\
    er/sie/es & \textbf{bedeutete}\\
    wir & \textbf{bedeuteten}\\
    ihr & \textbf{bedeutetet}\\
    sie/Sie & \textbf{bedeuteten}\\
  \end{tabular}
  \hfill
  \begin{tabular}{rl}
     \multicolumn{2}{l}{\textbf{IMPERATIV} }\\
   & \textbf{bedeut(e) (du)!}\\
   & \textbf{bedeuten wir!}\\
   & \textbf{bedeutet (ihr)!}\\
   & \textbf{bedeuten Sie!}\\
   & \vspace{10pt}
  \end{tabular}

    \vspace{10pt}

 \begin{tabular}{rl}
     \multicolumn{2}{l}{\textbf{PERFEKT}}\\
   & haben + \textbf{bedeutet}\\
  \end{tabular}
\hfill
   \begin{tabular}{rl}
   
    \multicolumn{2}{l}{\textbf{FUTURE I} }\\
   & werden + \textbf{bedeuten}\\
 
  \end{tabular}
\hfill
   \begin{tabular}{rl}
      \multicolumn{2}{l}{\textbf{PLUSQUAMPERFEKT}}\\
   & hatten + \textbf{bedeutet}\\
  \end{tabular}
  
\vspace{10pt}

\begin{tabular}{lll}
\multicolumn{3}{l}{\textbf{MIT PRÄPOSITIONEN}}\\
& \textbf{---} & \\
\end{tabular}
\hfill
\begin{tabular}{lll}
\multicolumn{3}{l}{\textbf{TRENNBARE VERBEN}}\\
& \textbf{---} & \\
\end{tabular}

\vspace{-5pt}
\end{verbentry}

% beginnen (to begin)
\begin{verbentry}{
  style = {acc},
  toc = {beginnen (to begin)},
  name = {beginnen},
  english = {to begin},
  type = {irregular, + Akkusativ},
  auxiliary = {haben}
}
 
    
     \vspace{5pt}

    \hrule
    
    \vspace{5pt}

  \begin{tabular}{rl}
     \multicolumn{2}{l}{\textbf{PRÄSENS} }\\
    ich & \textbf{beginne}\\
    du & \textbf{beginnst}\\
    er/sie/es & \textbf{beginnt}\\
    wir & \textbf{beginnen}\\
    ihr & \textbf{beginnt}\\
    sie/Sie & \textbf{beginnen}\\
  \end{tabular}
  \hfill
     \begin{tabular}{rl}
    \multicolumn{2}{l}{\textbf{PRÄTERITUM} }\\
    ich & \textbf{begann}\\
    du & \textbf{begannst}\\
    er/sie/es & \textbf{begann}\\
    wir & \textbf{begannen}\\
    ihr & \textbf{begannt}\\
    sie/Sie & \textbf{begannen}\\
  \end{tabular}
  \hfill
  \begin{tabular}{rl}
     \multicolumn{2}{l}{\textbf{IMPERATIV} }\\
   & \textbf{beginn(e) (du)!}\\
   & \textbf{beginnen wir!}\\
   & \textbf{beginnt (ihr)!}\\
   & \textbf{beginnen Sie!}\\
   & \vspace{10pt}
  \end{tabular}

    \vspace{10pt}

 \begin{tabular}{rl}
     \multicolumn{2}{l}{\textbf{PERFEKT}}\\
  &  haben + \textbf{begonnen}\\
  \end{tabular}
\hfill
   \begin{tabular}{rl}
   
    \multicolumn{2}{l}{\textbf{FUTURE I} }\\
  &  werden + \textbf{beginnen}\\
 
  \end{tabular}
\hfill
   \begin{tabular}{rl}
      \multicolumn{2}{l}{\textbf{PLUSQUAMPERFEKT}}\\
   & hatten + \textbf{begonnen}\\
  \end{tabular}
  
\vspace{10pt}

\begin{tabular}{lll}
\multicolumn{3}{l}{\textbf{MIT PRÄPOSITIONEN}}\\
& \textbf{---} & \\
\end{tabular}
\hfill
\begin{tabular}{lll}
\multicolumn{3}{l}{\textbf{TRENNBARE VERBEN}}\\
& \textbf{---} & \\
\end{tabular}

\vspace{-5pt}
\end{verbentry}

% begrüßen (to greet)
\begin{verbentry}{
  style = {acc},
  toc = {begrüßen (to greet)},
  name = {begrüßen},
  english = {to greet},
  type = {regular, + Akkusativ},
  auxiliary = {haben}
}


\vspace{5pt}
\hrule
\vspace{5pt}

\begin{tabular}{rl}
\multicolumn{2}{l}{\textbf{PRÄSENS}}\\
ich & \textbf{begrüße}\\
du & \textbf{begrüßt}\\
er/sie/es & \textbf{begrüßt}\\
wir & \textbf{begrüßen}\\
ihr & \textbf{begrüßt}\\
sie/Sie & \textbf{begrüßen}\\
\end{tabular}
\hfill
\begin{tabular}{rl}
\multicolumn{2}{l}{\textbf{PRÄTERITUM}}\\
ich & \textbf{begrüßte}\\
du & \textbf{begrüßtest}\\
er/sie/es & \textbf{begrüßte}\\
wir & \textbf{begrüßten}\\
ihr & \textbf{begrüßtet}\\
sie/Sie & \textbf{begrüßten}\\
\end{tabular}
\hfill
\begin{tabular}{rl}
\multicolumn{2}{l}{\textbf{IMPERATIV}}\\
& \textbf{begrüße (du)!}\\
& \textbf{begrüßen wir!}\\
& \textbf{begrüßt (ihr)!}\\
& \textbf{begrüßen Sie!}\\
\end{tabular}

\vspace{10pt}

\begin{tabular}{rl}
\multicolumn{2}{l}{\textbf{PERFEKT}}\\
& haben + \textbf{begrüßt}\\
\end{tabular}
\hfill
\begin{tabular}{rl}
\multicolumn{2}{l}{\textbf{FUTURE I}}\\
& werden + \textbf{begrüßen}\\
\end{tabular}
\hfill
\begin{tabular}{rl}
\multicolumn{2}{l}{\textbf{PLUSQUAMPERFEKT}}\\
& hatten + \textbf{begrüßt}\\
\end{tabular}

\vspace{10pt}

\begin{tabular}{lll}
\multicolumn{3}{l}{\textbf{MIT PRÄPOSITIONEN}}\\
& \textbf{---} & \\
\end{tabular}
\hfill
\begin{tabular}{lll}
\multicolumn{3}{l}{\textbf{TRENNBARE VERBEN}}\\
& \textbf{---} & \\
\end{tabular}

\vspace{-5pt}
\end{verbentry}

% behandeln (to treat)
\begin{verbentry}{
  style = {default},
  toc = {behandeln (to treat)},
  name = {behandeln},
  english = {to treat},
  type = {regular},
  auxiliary = {haben}
}
 
    
     \vspace{5pt}

    \hrule
    
    \vspace{5pt}

  \begin{tabular}{rl}
     \multicolumn{2}{l}{\textbf{PRÄSENS} }\\
    ich & \textbf{behandle}\\
    du & \textbf{behandelst}\\
    er/sie/es & \textbf{behandelt}\\
    wir & \textbf{behandeln}\\
    ihr & \textbf{behandelt}\\
    sie/Sie & \textbf{behandeln}\\
  \end{tabular}
  \hfill
     \begin{tabular}{rl}
    \multicolumn{2}{l}{\textbf{PRÄTERITUM} }\\
    ich & \textbf{behandelte}\\
    du & \textbf{behandeltest}\\
    er/sie/es & \textbf{behandelte}\\
    wir & \textbf{behandelten}\\
    ihr & \textbf{behandeltet}\\
    sie/Sie & \textbf{behandelten}\\
  \end{tabular}
  \hfill
  \begin{tabular}{rl}
     \multicolumn{2}{l}{\textbf{IMPERATIV} }\\
   & \textbf{behandle (du)!}\\
   & \textbf{behandeln wir!}\\
   & \textbf{behandelt (ihr)!}\\
   & \textbf{behandeln Sie!}\\
   & \vspace{10pt}
  \end{tabular}

    \vspace{10pt}

 \begin{tabular}{rl}
     \multicolumn{2}{l}{\textbf{PERFEKT}}\\
   & haben + \textbf{behandelt}\\
  \end{tabular}
\hfill
   \begin{tabular}{rl}
    \multicolumn{2}{l}{\textbf{FUTURE I} }\\
   & werde + \textbf{behandeln}\\
  \end{tabular}
\hfill
   \begin{tabular}{rl}
      \multicolumn{2}{l}{\textbf{PLUSQUAMPERFEKT}}\\
   & hatten + \textbf{behandelt}\\
  \end{tabular}

   \vspace{10pt}

   \begin{tabular}{lll}
      \multicolumn{3}{l}{\textbf{MIT PRÄPOSITIONEN}}\\
& \textbf{behandeln mit} + Dat & (to treat with)\\
& \textbf{behandeln wegen} + Gen & (to treat because of / for)\\
\end{tabular}
  \hfill
   \begin{tabular}{lll}
\multicolumn{3}{l}{\textbf{TRENNBARE VERBEN}}\\
& \textbf{---} & \\
\end{tabular}

  \vspace{-5pt}
\end{verbentry}

% beißen (to bite)
\begin{verbentry}{
  style = {default},
  toc = {beißen (to bite)},
  name = {beißen},
  english = {to bite},
  type = {irregular, strong},
  auxiliary = {haben}
}
 
    
     \vspace{5pt}

    \hrule
    
    \vspace{5pt}

  \begin{tabular}{rl}
     \multicolumn{2}{l}{\textbf{PRÄSENS} }\\
    ich & \textbf{beiße}\\
    du & \textbf{beißt}\\
    er/sie/es & \textbf{beißt}\\
    wir & \textbf{beißen}\\
    ihr & \textbf{beißt}\\
    sie/Sie & \textbf{beißen}\\
  \end{tabular}
  \hfill
     \begin{tabular}{rl}
    \multicolumn{2}{l}{\textbf{PRÄTERITUM} }\\
    ich & \textbf{biss}\\
    du & \textbf{bissest}\\
    er/sie/es & \textbf{biss}\\
    wir & \textbf{bissen}\\
    ihr & \textbf{bisst}\\
    sie/Sie & \textbf{bissen}\\
  \end{tabular}
  \hfill
  \begin{tabular}{rl}
     \multicolumn{2}{l}{\textbf{IMPERATIV} }\\
   & \textbf{beiß(e) (du)!}\\
   & \textbf{beißen wir!}\\
   & \textbf{beißt (ihr)!}\\
   & \textbf{beißen Sie!}\\
   & \vspace{10pt}
  \end{tabular}

    \vspace{10pt}

 \begin{tabular}{rl}
     \multicolumn{2}{l}{\textbf{PERFEKT}}\\
   & haben + \textbf{gebissen}\\
  \end{tabular}
\hfill
   \begin{tabular}{rl}
    \multicolumn{2}{l}{\textbf{FUTURE I} }\\
   & werde + \textbf{beißen}\\
  \end{tabular}
\hfill
   \begin{tabular}{rl}
      \multicolumn{2}{l}{\textbf{PLUSQUAMPERFEKT}}\\
   & hatten + \textbf{gebissen}\\
  \end{tabular}

   \vspace{10pt}

   \begin{tabular}{lll}
      \multicolumn{3}{l}{\textbf{MIT PRÄPOSITIONEN}}\\
& \textbf{beißen in} + Akk & (to bite into)\\
& \textbf{beißen nach} + Dat & (to snap at / bite at)\\
\end{tabular}
  \hfill
   \begin{tabular}{lll}
      \multicolumn{3}{l}{\textbf{TRENNBARE VERBEN}}\\
 & \textbf{ab$|$beißen} & (to bite off)\\
 & \textbf{zu$|$beißen} & (to bite down / snap)\\
\vspace{10pt}
  \end{tabular}

  \vspace{-5pt}
\end{verbentry}

% bekommen (to receive)
\begin{verbentry}{
  style = {acc},
  toc = {bekommen (to receive)},
  name = {bekommen},
  english = {to receive},
  type = {irregular, + Akkusativ},
  auxiliary = {haben}
}
 
    
     \vspace{5pt}

    \hrule
    
    \vspace{5pt}

  \begin{tabular}{rl}
     \multicolumn{2}{l}{\textbf{PRÄSENS} }\\
    ich & \textbf{bekomme}\\
    du & \textbf{bekommst}\\
    er/sie/es & \textbf{bekommt}\\
    wir & \textbf{bekommen}\\
    ihr & \textbf{bekommt}\\
    sie/Sie & \textbf{bekommen}\\
  \end{tabular}
  \hfill
     \begin{tabular}{rl}
    \multicolumn{2}{l}{\textbf{PRÄTERITUM} }\\
    ich & \textbf{bekam}\\
    du & \textbf{bekamst}\\
    er/sie/es & \textbf{bekam}\\
    wir & \textbf{bekamen}\\
    ihr & \textbf{bekamt}\\
    sie/Sie & \textbf{bekamen}\\
  \end{tabular}
  \hfill
  \begin{tabular}{rl}
     \multicolumn{2}{l}{\textbf{IMPERATIV} }\\
   & \textbf{bekomm(e) (du)!}\\
   & \textbf{bekommen wir!}\\
   & \textbf{bekommt (ihr)!}\\
   & \textbf{bekommen Sie!}\\
   & \vspace{10pt}
  \end{tabular}

    \vspace{10pt}

 \begin{tabular}{rl}
     \multicolumn{2}{l}{\textbf{PERFEKT}}\\
  &  haben + \textbf{bekommen}\\
  \end{tabular}
\hfill
   \begin{tabular}{rl}
   
    \multicolumn{2}{l}{\textbf{FUTURE I} }\\
  &  werden + \textbf{bekommen}\\
 
  \end{tabular}
\hfill
   \begin{tabular}{rl}
      \multicolumn{2}{l}{\textbf{PLUSQUAMPERFEKT}}\\
  &  hatten + \textbf{bekommen}\\
  \end{tabular}
  
\vspace{10pt}

\begin{tabular}{lll}
\multicolumn{3}{l}{\textbf{MIT PRÄPOSITIONEN}}\\
& \textbf{---} & \\
\end{tabular}
\hfill
\begin{tabular}{lll}
\multicolumn{3}{l}{\textbf{TRENNBARE VERBEN}}\\
& \textbf{---} & \\
\end{tabular}

\vspace{-5pt}
\end{verbentry}

% belasten (to burden / to strain)
\begin{verbentry}{
  style = {default},
  toc = {belasten (to burden / to strain)},
  name = {belasten},
  english = {to burden / to strain},
  type = {regular},
  auxiliary = {haben}
}
 
    
     \vspace{5pt}

    \hrule
    
    \vspace{5pt}

  \begin{tabular}{rl}
     \multicolumn{2}{l}{\textbf{PRÄSENS} }\\
    ich & \textbf{belaste}\\
    du & \textbf{belastest}\\
    er/sie/es & \textbf{belastet}\\
    wir & \textbf{belasten}\\
    ihr & \textbf{belastet}\\
    sie/Sie & \textbf{belasten}\\
  \end{tabular}
  \hfill
     \begin{tabular}{rl}
    \multicolumn{2}{l}{\textbf{PRÄTERITUM} }\\
    ich & \textbf{belastete}\\
    du & \textbf{belastetest}\\
    er/sie/es & \textbf{belastete}\\
    wir & \textbf{belasteten}\\
    ihr & \textbf{belastetet}\\
    sie/Sie & \textbf{belasteten}\\
  \end{tabular}
  \hfill
  \begin{tabular}{rl}
     \multicolumn{2}{l}{\textbf{IMPERATIV} }\\
   & \textbf{belaste (du)!}\\
   & \textbf{belasten wir!}\\
   & \textbf{belastet (ihr)!}\\
   & \textbf{belasten Sie!}\\
   & \vspace{10pt}
  \end{tabular}

    \vspace{10pt}

 \begin{tabular}{rl}
     \multicolumn{2}{l}{\textbf{PERFEKT}}\\
   & haben + \textbf{belastet}\\
  \end{tabular}
\hfill
   \begin{tabular}{rl}
    \multicolumn{2}{l}{\textbf{FUTURE I} }\\
   & werde + \textbf{belasten}\\
  \end{tabular}
\hfill
   \begin{tabular}{rl}
      \multicolumn{2}{l}{\textbf{PLUSQUAMPERFEKT}}\\
   & hatten + \textbf{belastet}\\
  \end{tabular}

   \vspace{10pt}

   \begin{tabular}{lll}
      \multicolumn{3}{l}{\textbf{MIT PRÄPOSITIONEN}}\\
& \textbf{belasten mit} + Dat & (to burden with)\\
& \textbf{belasten durch} + Akk & (to strain/burden through)\\
\end{tabular}
  \hfill
   \begin{tabular}{lll}
\multicolumn{3}{l}{\textbf{TRENNBARE VERBEN}}\\
& \textbf{---} & \\
\end{tabular}

  \vspace{-5pt}
\end{verbentry}

% belegen (to occupy / register / cover)
\begin{verbentry}{
  style = {acc},
  toc = {belegen (to occupy / register / cover)},
  name = {belegen},
  english = {to occupy / register / cover},
  type = {regular, + Akkusativ},
  auxiliary = {haben}
}


\vspace{5pt}
\hrule
\vspace{5pt}

\begin{tabular}{rl}
\multicolumn{2}{l}{\textbf{PRÄSENS}}\\
ich & \textbf{belege}\\
du & \textbf{belegst}\\
er/sie/es & \textbf{belegt}\\
wir & \textbf{belegen}\\
ihr & \textbf{belegt}\\
sie/Sie & \textbf{belegen}\\
\end{tabular}
\hfill
\begin{tabular}{rl}
\multicolumn{2}{l}{\textbf{PRÄTERITUM}}\\
ich & \textbf{belegte}\\
du & \textbf{belegtest}\\
er/sie/es & \textbf{belegte}\\
wir & \textbf{belegten}\\
ihr & \textbf{belegtet}\\
sie/Sie & \textbf{belegten}\\
\end{tabular}
\hfill
\begin{tabular}{rl}
\multicolumn{2}{l}{\textbf{IMPERATIV}}\\
& \textbf{belege (du)!}\\
& \textbf{belegen wir!}\\
& \textbf{belegt (ihr)!}\\
& \textbf{belegen Sie!}\\
\end{tabular}

\vspace{10pt}

\begin{tabular}{rl}
\multicolumn{2}{l}{\textbf{PERFEKT}}\\
& haben + \textbf{belegt}\\
\end{tabular}
\hfill
\begin{tabular}{rl}
\multicolumn{2}{l}{\textbf{FUTURE I}}\\
& werden + \textbf{belegen}\\
\end{tabular}
\hfill
\begin{tabular}{rl}
\multicolumn{2}{l}{\textbf{PLUSQUAMPERFEKT}}\\
& hatten + \textbf{belegt}\\
\end{tabular}

\vspace{10pt}

\begin{tabular}{lll}
\multicolumn{3}{l}{\textbf{MIT PRÄPOSITIONEN}}\\
& \textbf{belegen mit} + Dat & (to cover with)\\
\end{tabular}
\hfill
\begin{tabular}{lll}
\multicolumn{3}{l}{\textbf{TRENNBARE VERBEN}}\\
& \textbf{---} & \\
\end{tabular}
\vspace{-5pt}
\end{verbentry}

% bemerken (to notice)
\begin{verbentry}{
  style = {acc},
  toc = {bemerken (to notice)},
  name = {bemerken},
  english = {to notice},
  type = {regular, + Akkusativ},
  auxiliary = {haben}
}


\vspace{5pt}
\hrule
\vspace{5pt}

\begin{tabular}{rl}
\multicolumn{2}{l}{\textbf{PRÄSENS}}\\
ich & \textbf{bemerke}\\
du & \textbf{bemerkst}\\
er/sie/es & \textbf{bemerkt}\\
wir & \textbf{bemerken}\\
ihr & \textbf{bemerkt}\\
sie/Sie & \textbf{bemerken}\\
\end{tabular}
\hfill
\begin{tabular}{rl}
\multicolumn{2}{l}{\textbf{PRÄTERITUM}}\\
ich & \textbf{bemerkte}\\
du & \textbf{bemerktest}\\
er/sie/es & \textbf{bemerkte}\\
wir & \textbf{bemerkten}\\
ihr & \textbf{bemerktet}\\
sie/Sie & \textbf{bemerkten}\\
\end{tabular}
\hfill
\begin{tabular}{rl}
\multicolumn{2}{l}{\textbf{IMPERATIV}}\\
& \textbf{bemerke (du)!}\\
& \textbf{bemerken wir!}\\
& \textbf{bemerkt (ihr)!}\\
& \textbf{bemerken Sie!}\\
\end{tabular}

\vspace{10pt}

\begin{tabular}{rl}
\multicolumn{2}{l}{\textbf{PERFEKT}}\\
& haben + \textbf{bemerkt}\\
\end{tabular}
\hfill
\begin{tabular}{rl}
\multicolumn{2}{l}{\textbf{FUTURE I}}\\
& werden + \textbf{bemerken}\\
\end{tabular}
\hfill
\begin{tabular}{rl}
\multicolumn{2}{l}{\textbf{PLUSQUAMPERFEKT}}\\
& hatten + \textbf{bemerkt}\\
\end{tabular}

\vspace{10pt}

\begin{tabular}{lll}
\multicolumn{3}{l}{\textbf{MIT PRÄPOSITIONEN}}\\
& \textbf{---} & \\
\end{tabular}
\hfill
\begin{tabular}{lll}
\multicolumn{3}{l}{\textbf{TRENNBARE VERBEN}}\\
& \textbf{---} & \\
\end{tabular}

\vspace{-5pt}
\end{verbentry}

% benutzen (to use)
\begin{verbentry}{
  style = {acc},
  toc = {benutzen (to use)},
  name = {benutzen},
  english = {to use},
  type = {regular, + Akkusativ},
  auxiliary = {haben}
}
 
    
     \vspace{5pt}

    \hrule
    
    \vspace{5pt}

  \begin{tabular}{rl}
     \multicolumn{2}{l}{\textbf{PRÄSENS} }\\
    ich & \textbf{benutze}\\
    du & \textbf{benutzt}\\
    er/sie/es & \textbf{benutz}\\
    wir & \textbf{benutzen}\\
    ihr & \textbf{benutzt}\\
    sie/Sie & \textbf{benutzen}\\
  \end{tabular}
  \hfill
     \begin{tabular}{rl}
    \multicolumn{2}{l}{\textbf{PRÄTERITUM} }\\
    ich & \textbf{benutzte}\\
    du & \textbf{benutztest}\\
    er/sie/es & \textbf{benutzte}\\
    wir & \textbf{benutzten}\\
    ihr & \textbf{benutztet}\\
    sie/Sie & \textbf{benutzten}\\
  \end{tabular}
  \hfill
  \begin{tabular}{rl}
     \multicolumn{2}{l}{\textbf{IMPERATIV} }\\
   & \textbf{benutz(e) (du)!}\\
   & \textbf{benutzen wir!}\\
   & \textbf{benutzt (ihr)!}\\
   & \textbf{benutzen Sie!}\\
   & \vspace{10pt}
  \end{tabular}

    \vspace{10pt}

 \begin{tabular}{rl}
     \multicolumn{2}{l}{\textbf{PERFEKT}}\\
   & haben + \textbf{benutzt}\\
  \end{tabular}
\hfill
   \begin{tabular}{rl}
   
    \multicolumn{2}{l}{\textbf{FUTURE I} }\\
   & werden + \textbf{benutzen}\\
 
  \end{tabular}
\hfill
   \begin{tabular}{rl}
      \multicolumn{2}{l}{\textbf{PLUSQUAMPERFEKT}}\\
   & hatten + \textbf{benutzt}\\
  \end{tabular}
  
  \vspace{-5pt}\vspace{10pt}

\begin{tabular}{lll}
\multicolumn{3}{l}{\textbf{MIT PRÄPOSITIONEN}}\\
& \textbf{---} & \\
\end{tabular}
\hfill
\begin{tabular}{lll}
\multicolumn{3}{l}{\textbf{TRENNBARE VERBEN}}\\
& \textbf{---} & \\
\end{tabular}

\vspace{-5pt}
\end{verbentry}

% beobachten (to observe)
\begin{verbentry}{
  style = {acc},
  toc = {beobachten (to observe)},
  name = {beobachten},
  english = {to observe},
  type = {regular, + Akkusativ},
  auxiliary = {haben}
}
 

\vspace{5pt}
\hrule
\vspace{5pt}

\begin{tabular}{rl}
\multicolumn{2}{l}{\textbf{PRÄSENS}}\\
ich & \textbf{beobachte}\\
du & \textbf{beobachtest}\\
er/sie/es & \textbf{beobachtet}\\
wir & \textbf{beobachten}\\
ihr & \textbf{beobachtet}\\
sie/Sie & \textbf{beobachten}\\
\end{tabular}
\hfill
\begin{tabular}{rl}
\multicolumn{2}{l}{\textbf{PRÄTERITUM}}\\
ich & \textbf{beobachtete}\\
du & \textbf{beobachtetest}\\
er/sie/es & \textbf{beobachtete}\\
wir & \textbf{beobachteten}\\
ihr & \textbf{beobachtetet}\\
sie/Sie & \textbf{beobachteten}\\
\end{tabular}
\hfill
\begin{tabular}{rl}
\multicolumn{2}{l}{\textbf{IMPERATIV}}\\
& \textbf{beobachte (du)!}\\
& \textbf{beobachten wir!}\\
& \textbf{beobachtet (ihr)!}\\
& \textbf{beobachten Sie!}\\
\end{tabular}

\vspace{10pt}

\begin{tabular}{rl}
\multicolumn{2}{l}{\textbf{PERFEKT}}\\
& haben + \textbf{beobachtet}\\
\end{tabular}
\hfill
\begin{tabular}{rl}
\multicolumn{2}{l}{\textbf{FUTURE I}}\\
& werden + \textbf{beobachten}\\
\end{tabular}
\hfill
\begin{tabular}{rl}
\multicolumn{2}{l}{\textbf{PLUSQUAMPERFEKT}}\\
& hatten + \textbf{beobachtet}\\
\end{tabular}

\vspace{10pt}

\begin{tabular}{lll}
\multicolumn{3}{l}{\textbf{MIT PRÄPOSITIONEN}}\\
& \textbf{beobachten Akk} & (to observe something)\\
\end{tabular}
\hfill
\begin{tabular}{lll}
\multicolumn{3}{l}{\textbf{TRENNBARE VERBEN}}\\
& \textbf{---} & \\
\end{tabular}

\vspace{-5pt}
\end{verbentry}

% beraten (to advise)
\begin{verbentry}{
  style = {acc},
  toc = {beraten (to advise)},
  name = {beraten},
  english = {to advise},
  type = {irregular, + Akkusativ},
  auxiliary = {haben}
}


\vspace{5pt}
\hrule
\vspace{5pt}

\begin{tabular}{rl}
\multicolumn{2}{l}{\textbf{PRÄSENS}}\\
ich & \textbf{berate}\\
du & \textbf{berätst}\\
er/sie/es & \textbf{berät}\\
wir & \textbf{beraten}\\
ihr & \textbf{beratet}\\
sie/Sie & \textbf{beraten}\\
\end{tabular}
\hfill
\begin{tabular}{rl}
\multicolumn{2}{l}{\textbf{PRÄTERITUM}}\\
ich & \textbf{beriet}\\
du & \textbf{berietst}\\
er/sie/es & \textbf{beriet}\\
wir & \textbf{berieten}\\
ihr & \textbf{berietet}\\
sie/Sie & \textbf{berieten}\\
\end{tabular}
\hfill
\begin{tabular}{rl}
\multicolumn{2}{l}{\textbf{IMPERATIV}}\\
& \textbf{berate (du)!}\\
& \textbf{beraten wir!}\\
& \textbf{beratet (ihr)!}\\
& \textbf{beraten Sie!}\\
\end{tabular}

\vspace{10pt}

\begin{tabular}{rl}
\multicolumn{2}{l}{\textbf{PERFEKT}}\\
& haben + \textbf{beraten}\\
\end{tabular}
\hfill
\begin{tabular}{rl}
\multicolumn{2}{l}{\textbf{FUTURE I}}\\
& werden + \textbf{beraten}\\
\end{tabular}
\hfill
\begin{tabular}{rl}
\multicolumn{2}{l}{\textbf{PLUSQUAMPERFEKT}}\\
& hatten + \textbf{beraten}\\
\end{tabular}

\vspace{10pt}

\begin{tabular}{lll}
\multicolumn{3}{l}{\textbf{MIT PRÄPOSITIONEN}}\\
& \textbf{beraten über} + Akk & (to advise about)\\
& \textbf{beraten bei} + Dat & (to advise during / assist with)\\
\end{tabular}
\hfill
\begin{tabular}{lll}
\multicolumn{3}{l}{\textbf{TRENNBARE VERBEN}}\\
& \textbf{---} & \\
\end{tabular}

\vspace{-5pt}
\end{verbentry}

% bereiten (to prepare / cause)
\begin{verbentry}{
  style = {acc},
  toc = {bereiten (to prepare / cause)},
  name = {bereiten},
  english = {to prepare / cause},
  type = {regular, + Akkusativ},
  auxiliary = {haben}
}


\vspace{5pt}
\hrule
\vspace{5pt}

\begin{tabular}{rl}
\multicolumn{2}{l}{\textbf{PRÄSENS}}\\
ich & \textbf{bereite}\\
du & \textbf{bereitest}\\
er/sie/es & \textbf{bereitet}\\
wir & \textbf{bereiten}\\
ihr & \textbf{bereitet}\\
sie/Sie & \textbf{bereiten}\\
\end{tabular}
\hfill
\begin{tabular}{rl}
\multicolumn{2}{l}{\textbf{PRÄTERITUM}}\\
ich & \textbf{bereitete}\\
du & \textbf{bereitetest}\\
er/sie/es & \textbf{bereitete}\\
wir & \textbf{bereiteten}\\
ihr & \textbf{bereitetet}\\
sie/Sie & \textbf{bereiteten}\\
\end{tabular}
\hfill
\begin{tabular}{rl}
\multicolumn{2}{l}{\textbf{IMPERATIV}}\\
& \textbf{bereite (du)!}\\
& \textbf{bereiten wir!}\\
& \textbf{bereitet (ihr)!}\\
& \textbf{bereiten Sie!}\\
\end{tabular}

\vspace{10pt}

\begin{tabular}{rl}
\multicolumn{2}{l}{\textbf{PERFEKT}}\\
& haben + \textbf{bereitet}\\
\end{tabular}
\hfill
\begin{tabular}{rl}
\multicolumn{2}{l}{\textbf{FUTURE I}}\\
& werden + \textbf{bereiten}\\
\end{tabular}
\hfill
\begin{tabular}{rl}
\multicolumn{2}{l}{\textbf{PLUSQUAMPERFEKT}}\\
& hatten + \textbf{bereitet}\\
\end{tabular}

\vspace{10pt}

\begin{tabular}{lll}
\multicolumn{3}{l}{\textbf{MIT PRÄPOSITIONEN}}\\
& \textbf{bereiten auf} + Akk & (to prepare for)\\
\end{tabular}
\hfill
\begin{tabular}{lll}
\multicolumn{3}{l}{\textbf{TRENNBARE VERBEN}}\\
& \textbf{vor$|$bereiten} & (to prepare)\\
\end{tabular}

\vspace{-5pt}
\end{verbentry}

% besetzen (to occupy / fill a position)
\begin{verbentry}{
  style = {acc},
  toc = {besetzen (to occupy / fill a position)},
  name = {besetzen},
  english = {to occupy / fill a position},
  type = {regular, + Akkusativ},
  auxiliary = {haben}
}


\vspace{5pt}
\hrule
\vspace{5pt}

\begin{tabular}{rl}
\multicolumn{2}{l}{\textbf{PRÄSENS}}\\
ich & \textbf{besetze}\\
du & \textbf{besetzt}\\
er/sie/es & \textbf{besetzt}\\
wir & \textbf{besetzen}\\
ihr & \textbf{besetzt}\\
sie/Sie & \textbf{besetzen}\\
\end{tabular}
\hfill
\begin{tabular}{rl}
\multicolumn{2}{l}{\textbf{PRÄTERITUM}}\\
ich & \textbf{besetzte}\\
du & \textbf{besetztest}\\
er/sie/es & \textbf{besetzte}\\
wir & \textbf{besetzten}\\
ihr & \textbf{besetztet}\\
sie/Sie & \textbf{besetzten}\\
\end{tabular}
\hfill
\begin{tabular}{rl}
\multicolumn{2}{l}{\textbf{IMPERATIV}}\\
& \textbf{besetze (du)!}\\
& \textbf{besetzen wir!}\\
& \textbf{besetzt (ihr)!}\\
& \textbf{besetzen Sie!}\\
\end{tabular}

\vspace{10pt}

\begin{tabular}{rl}
\multicolumn{2}{l}{\textbf{PERFEKT}}\\
& haben + \textbf{besetzt}\\
\end{tabular}
\hfill
\begin{tabular}{rl}
\multicolumn{2}{l}{\textbf{FUTURE I}}\\
& werden + \textbf{besetzen}\\
\end{tabular}
\hfill
\begin{tabular}{rl}
\multicolumn{2}{l}{\textbf{PLUSQUAMPERFEKT}}\\
& hatten + \textbf{besetzt}\\
\end{tabular}

\vspace{10pt}

\begin{tabular}{lll}
\multicolumn{3}{l}{\textbf{MIT PRÄPOSITIONEN}}\\
& \textbf{besetzen mit} + Dat & (to fill with / occupy with)\\
\end{tabular}
\hfill
\begin{tabular}{lll}
\multicolumn{3}{l}{\textbf{TRENNBARE VERBEN}}\\
& \textbf{---} & \\
\end{tabular}
\vspace{-5pt}
\end{verbentry}

% beschließen (to decide)
\begin{verbentry}{
  style = {acc},
  toc = {beschließen (to decide)},
  name = {beschließen},
  english = {to decide},
  type = {irregular, + Akkusativ},
  auxiliary = {haben}
}


\vspace{5pt}
\hrule
\vspace{5pt}

\begin{tabular}{rl}
\multicolumn{2}{l}{\textbf{PRÄSENS}}\\
ich & \textbf{beschließe}\\
du & \textbf{beschließt}\\
er/sie/es & \textbf{beschließt}\\
wir & \textbf{beschließen}\\
ihr & \textbf{beschließt}\\
sie/Sie & \textbf{beschließen}\\
\end{tabular}
\hfill
\begin{tabular}{rl}
\multicolumn{2}{l}{\textbf{PRÄTERITUM}}\\
ich & \textbf{beschloss}\\
du & \textbf{beschlossest}\\
er/sie/es & \textbf{beschloss}\\
wir & \textbf{beschlossen}\\
ihr & \textbf{beschlosst}\\
sie/Sie & \textbf{beschlossen}\\
\end{tabular}
\hfill
\begin{tabular}{rl}
\multicolumn{2}{l}{\textbf{IMPERATIV}}\\
& \textbf{beschließ(e) (du)!}\\
& \textbf{beschließen wir!}\\
& \textbf{beschließt (ihr)!}\\
& \textbf{beschließen Sie!}\\
\end{tabular}

\vspace{10pt}

\begin{tabular}{rl}
\multicolumn{2}{l}{\textbf{PERFEKT}}\\
& haben + \textbf{beschlossen}\\
\end{tabular}
\hfill
\begin{tabular}{rl}
\multicolumn{2}{l}{\textbf{FUTURE I}}\\
& werden + \textbf{beschließen}\\
\end{tabular}
\hfill
\begin{tabular}{rl}
\multicolumn{2}{l}{\textbf{PLUSQUAMPERFEKT}}\\
& hatten + \textbf{beschlossen}\\
\end{tabular}

\vspace{10pt}

\begin{tabular}{lll}
\multicolumn{3}{l}{\textbf{MIT PRÄPOSITIONEN}}\\
& \textbf{beschließen Akk} & (to decide something)\\
\end{tabular}
\hfill
\begin{tabular}{lll}
\multicolumn{3}{l}{\textbf{TRENNBARE VERBEN}}\\
& \textbf{---} & \\
\end{tabular}

\vspace{-5pt}
\end{verbentry}

% beschreiben (to describe)
\begin{verbentry}{
  style = {acc},
  toc = {beschreiben (to describe)},
  name = {beschreiben},
  english = {to describe},
  type = {irregular, + Akkusativ},
  auxiliary = {haben}
}
 
    
     \vspace{5pt}

    \hrule
    
    \vspace{5pt}

  \begin{tabular}{rl}
     \multicolumn{2}{l}{\textbf{PRÄSENS} }\\
    ich & \textbf{beschreibe}\\
    du & \textbf{beschreibst}\\
    er/sie/es & \textbf{beschreibt}\\
    wir & \textbf{beschreiben}\\
    ihr & \textbf{beschreibt}\\
    sie/Sie & \textbf{beschreiben}\\
  \end{tabular}
  \hfill
     \begin{tabular}{rl}
    \multicolumn{2}{l}{\textbf{PRÄTERITUM} }\\
    ich & \textbf{beschrieb}\\
    du & \textbf{beschriebst}\\
    er/sie/es & \textbf{beschrieb}\\
    wir & \textbf{beschrieben}\\
    ihr & \textbf{beschriebt}\\
    sie/Sie & \textbf{beschrieben}\\
  \end{tabular}
  \hfill
  \begin{tabular}{rl}
     \multicolumn{2}{l}{\textbf{IMPERATIV} }\\
   & \textbf{beschreib(e) (du)!}\\
   & \textbf{beschreiben wir!}\\
   & \textbf{beschreibt (ihr)!}\\
   & \textbf{beschreiben Sie!}\\
   & \vspace{10pt}
  \end{tabular}

    \vspace{10pt}

 \begin{tabular}{rl}
     \multicolumn{2}{l}{\textbf{PERFEKT}}\\
   & haben + \textbf{beschrieben}\\
  \end{tabular}
\hfill
   \begin{tabular}{rl}
   
    \multicolumn{2}{l}{\textbf{FUTURE I} }\\
   & werden + \textbf{beschreiben}\\
 
  \end{tabular}
\hfill
   \begin{tabular}{rl}
      \multicolumn{2}{l}{\textbf{PLUSQUAMPERFEKT}}\\
   & hatten + \textbf{beschrieben}\\
  \end{tabular}
  
\vspace{10pt}

\begin{tabular}{lll}
\multicolumn{3}{l}{\textbf{MIT PRÄPOSITIONEN}}\\
& \textbf{---} & \\
\end{tabular}
\hfill
\begin{tabular}{lll}
\multicolumn{3}{l}{\textbf{TRENNBARE VERBEN}}\\
& \textbf{---} & \\
\end{tabular}

  \vspace{-5pt}
\end{verbentry}

% bestätigen (to confirm)
\begin{verbentry}{
  style = {acc},
  toc = {bestätigen (to confirm)},
  name = {bestätigen},
  english = {to confirm},
  type = {regular, + Akkusativ},
  auxiliary = {haben}
}
 

\vspace{5pt}
\hrule
\vspace{5pt}

\begin{tabular}{rl}
\multicolumn{2}{l}{\textbf{PRÄSENS}}\\
ich & \textbf{bestätige}\\
du & \textbf{bestätigst}\\
er/sie/es & \textbf{bestätigt}\\
wir & \textbf{bestätigen}\\
ihr & \textbf{bestätigt}\\
sie/Sie & \textbf{bestätigen}\\
\end{tabular}
\hfill
\begin{tabular}{rl}
\multicolumn{2}{l}{\textbf{PRÄTERITUM}}\\
ich & \textbf{bestätigte}\\
du & \textbf{bestätigtest}\\
er/sie/es & \textbf{bestätigte}\\
wir & \textbf{bestätigten}\\
ihr & \textbf{bestätigtet}\\
sie/Sie & \textbf{bestätigten}\\
\end{tabular}
\hfill
\begin{tabular}{rl}
\multicolumn{2}{l}{\textbf{IMPERATIV}}\\
& \textbf{bestätige (du)!}\\
& \textbf{bestätigen wir!}\\
& \textbf{bestätigt (ihr)!}\\
& \textbf{bestätigen Sie!}\\
\end{tabular}

\vspace{10pt}

\begin{tabular}{rl}
\multicolumn{2}{l}{\textbf{PERFEKT}}\\
& haben + \textbf{bestätigt}\\
\end{tabular}
\hfill
\begin{tabular}{rl}
\multicolumn{2}{l}{\textbf{FUTURE I}}\\
& werden + \textbf{bestätigen}\\
\end{tabular}
\hfill
\begin{tabular}{rl}
\multicolumn{2}{l}{\textbf{PLUSQUAMPERFEKT}}\\
& hatten + \textbf{bestätigt}\\
\end{tabular}

\vspace{10pt}

\begin{tabular}{lll}
\multicolumn{3}{l}{\textbf{MIT PRÄPOSITIONEN}}\\
& \textbf{bestätigen gegenüber} + Dat & (to confirm to someone)\\
\end{tabular}
\hfill
\begin{tabular}{lll}
\multicolumn{3}{l}{\textbf{TRENNBARE VERBEN}}\\
& \textbf{---} & \\
\end{tabular}

\vspace{-5pt}
\end{verbentry}

% bestehen (to pass)
\begin{verbentry}{
  style = {acc},
  toc = {bestehen (to pass)},
  name = {bestehen},
  english = {to pass},
  type = {irregular, + Akkusativ},
  auxiliary = {haben}
}
 
    
     \vspace{5pt}

    \hrule
    
    \vspace{5pt}

  \begin{tabular}{rl}
     \multicolumn{2}{l}{\textbf{PRÄSENS} }\\
    ich & \textbf{bestehe}\\
    du & \textbf{bestehst}\\
    er/sie/es & \textbf{besteht}\\
    wir & \textbf{bestehen}\\
    ihr & \textbf{besteht}\\
    sie/Sie & \textbf{bestehen}\\
  \end{tabular}
  \hfill
     \begin{tabular}{rl}
    \multicolumn{2}{l}{\textbf{PRÄTERITUM} }\\
    ich & \textbf{bestand}\\
    du & \textbf{bestandest}\\
    er/sie/es & \textbf{bestand}\\
    wir & \textbf{bestanden}\\
    ihr & \textbf{bestandet}\\
    sie/Sie & \textbf{bestanden}\\
  \end{tabular}
  \hfill
  \begin{tabular}{rl}
     \multicolumn{2}{l}{\textbf{IMPERATIV} }\\
   & \textbf{besteh(e) (du)!}\\
   & \textbf{bestehen wir!}\\
   & \textbf{besteht (ihr)!}\\
   & \textbf{bestehen Sie!}\\
   & \vspace{10pt}
  \end{tabular}

    \vspace{10pt}

 \begin{tabular}{rl}
     \multicolumn{2}{l}{\textbf{PERFEKT}}\\
  &  haben + \textbf{bestanden}\\
  \end{tabular}
\hfill
   \begin{tabular}{rl}
   
    \multicolumn{2}{l}{\textbf{FUTURE I} }\\
  &  werden + \textbf{bestehen}\\
 
  \end{tabular}
\hfill
   \begin{tabular}{rl}
      \multicolumn{2}{l}{\textbf{PLUSQUAMPERFEKT}}\\
   & hatten + \textbf{bestanden}\\
  \end{tabular}
  
 \vspace{10pt}

\begin{tabular}{lll}
\multicolumn{3}{l}{\textbf{MIT PRÄPOSITIONEN}}\\
& \textbf{---} & \\
\end{tabular}
\hfill
\begin{tabular}{lll}
\multicolumn{3}{l}{\textbf{TRENNBARE VERBEN}}\\
& \textbf{---} & \\
\end{tabular}

\vspace{-5pt}
\end{verbentry}

% bestellen (to order)
\begin{verbentry}{
  style = {acc},
  toc = {bestellen (to order)},
  name = {bestellen},
  english = {to order},
  type = {regular, + Akkusativ},
  auxiliary = {haben}
}
 
    
     \vspace{5pt}

    \hrule
    
    \vspace{5pt}

  \begin{tabular}{rl}
     \multicolumn{2}{l}{\textbf{PRÄSENS} }\\
    ich & \textbf{bestelle}\\
    du & \textbf{bestellst}\\
    er/sie/es & \textbf{bestellt}\\
    wir & \textbf{bestellen}\\
    ihr & \textbf{bestellt}\\
    sie/Sie & \textbf{bestellen}\\
  \end{tabular}
  \hfill
     \begin{tabular}{rl}
    \multicolumn{2}{l}{\textbf{PRÄTERITUM} }\\
    ich & \textbf{bestellte}\\
    du & \textbf{bestelltest}\\
    er/sie/es & \textbf{bestellte}\\
    wir & \textbf{bestellten}\\
    ihr & \textbf{bestelltet}\\
    sie/Sie & \textbf{bestellten}\\
  \end{tabular}
  \hfill
  \begin{tabular}{rl}
     \multicolumn{2}{l}{\textbf{IMPERATIV} }\\
   & \textbf{bestell(e) (du)!}\\
   & \textbf{bestellen wir!}\\
   & \textbf{bestellt (ihr)!}\\
   & \textbf{bestellen Sie!}\\
   & \vspace{10pt}
  \end{tabular}

    \vspace{10pt}

 \begin{tabular}{rl}
     \multicolumn{2}{l}{\textbf{PERFEKT}}\\
   & haben + \textbf{bestellt}\\
  \end{tabular}
\hfill
   \begin{tabular}{rl}
   
    \multicolumn{2}{l}{\textbf{FUTURE I} }\\
   & werden + \textbf{bestellen}\\
 
  \end{tabular}
\hfill
   \begin{tabular}{rl}
      \multicolumn{2}{l}{\textbf{PLUSQUAMPERFEKT}}\\
   & hatten + \textbf{bestellt}\\
  \end{tabular}
  
 \vspace{10pt}

\begin{tabular}{lll}
\multicolumn{3}{l}{\textbf{MIT PRÄPOSITIONEN}}\\
& \textbf{---} & \\
\end{tabular}
\hfill
\begin{tabular}{lll}
\multicolumn{3}{l}{\textbf{TRENNBARE VERBEN}}\\
& \textbf{---} & \\
\end{tabular}

\vspace{-5pt}
\end{verbentry}

% betragen (to amount to / to be)
\begin{verbentry}{
  style = {acc},
  toc = {betragen (to amount to / to be)},
  name = {betragen},
  english = {to amount to / to be},
  type = {irregular, + Akkusativ},
  auxiliary = {haben}
}


\vspace{5pt}
\hrule
\vspace{5pt}

\begin{tabular}{rl}
\multicolumn{2}{l}{\textbf{PRÄSENS}}\\
ich & \textbf{betrage}\\
du & \textbf{beträgst}\\
er/sie/es & \textbf{beträgt}\\
wir & \textbf{betragen}\\
ihr & \textbf{betragt}\\
sie/Sie & \textbf{betragen}\\
\end{tabular}
\hfill
\begin{tabular}{rl}
\multicolumn{2}{l}{\textbf{PRÄTERITUM}}\\
ich & \textbf{betrug}\\
du & \textbf{betrugst}\\
er/sie/es & \textbf{betrug}\\
wir & \textbf{betrugen}\\
ihr & \textbf{betrugt}\\
sie/Sie & \textbf{betrugen}\\
\end{tabular}
\hfill
\begin{tabular}{rl}
\multicolumn{2}{l}{\textbf{IMPERATIV}}\\
& \textbf{---}\\
\end{tabular}

\vspace{10pt}

\begin{tabular}{rl}
\multicolumn{2}{l}{\textbf{PERFEKT}}\\
& haben + \textbf{betragen}\\
\end{tabular}
\hfill
\begin{tabular}{rl}
\multicolumn{2}{l}{\textbf{FUTURE I}}\\
& werden + \textbf{betragen}\\
\end{tabular}
\hfill
\begin{tabular}{rl}
\multicolumn{2}{l}{\textbf{PLUSQUAMPERFEKT}}\\
& hatten + \textbf{betragen}\\
\end{tabular}

\vspace{10pt}

\begin{tabular}{lll}
\multicolumn{3}{l}{\textbf{MIT PRÄPOSITIONEN}}\\
& \textbf{---} & \\
\end{tabular}
\hfill
\begin{tabular}{lll}
\multicolumn{3}{l}{\textbf{TRENNBARE VERBEN}}\\
& \textbf{---} & \\
\end{tabular}
\vspace{-5pt}
\end{verbentry}

% betreuen (to supervise / take care of)
\begin{verbentry}{
  style = {acc},
  toc = {betreuen (to supervise / take care of)},
  name = {betreuen},
  english = {to supervise / take care of},
  type = {regular, + Akkusativ},
  auxiliary = {haben}
}


\vspace{5pt}
\hrule
\vspace{5pt}

\begin{tabular}{rl}
\multicolumn{2}{l}{\textbf{PRÄSENS}}\\
ich & \textbf{betreue}\\
du & \textbf{betreust}\\
er/sie/es & \textbf{betreut}\\
wir & \textbf{betreuen}\\
ihr & \textbf{betreut}\\
sie/Sie & \textbf{betreuen}\\
\end{tabular}
\hfill
\begin{tabular}{rl}
\multicolumn{2}{l}{\textbf{PRÄTERITUM}}\\
ich & \textbf{betreute}\\
du & \textbf{betreutest}\\
er/sie/es & \textbf{betreute}\\
wir & \textbf{betreuten}\\
ihr & \textbf{betreutet}\\
sie/Sie & \textbf{betreuten}\\
\end{tabular}
\hfill
\begin{tabular}{rl}
\multicolumn{2}{l}{\textbf{IMPERATIV}}\\
& \textbf{betreue (du)!}\\
& \textbf{betreuen wir!}\\
& \textbf{betreut (ihr)!}\\
& \textbf{betreuen Sie!}\\
\end{tabular}

\vspace{10pt}

\begin{tabular}{rl}
\multicolumn{2}{l}{\textbf{PERFEKT}}\\
& haben + \textbf{betreut}\\
\end{tabular}
\hfill
\begin{tabular}{rl}
\multicolumn{2}{l}{\textbf{FUTURE I}}\\
& werden + \textbf{betreuen}\\
\end{tabular}
\hfill
\begin{tabular}{rl}
\multicolumn{2}{l}{\textbf{PLUSQUAMPERFEKT}}\\
& hatten + \textbf{betreut}\\
\end{tabular}

\vspace{10pt}

\begin{tabular}{lll}
\multicolumn{3}{l}{\textbf{MIT PRÄPOSITIONEN}}\\
& \textbf{betreuen bei} + Dat & (to assist during)\\
& \textbf{betreuen in} + Dat & (to supervise in)\\
\end{tabular}
\hfill
\begin{tabular}{lll}
\multicolumn{3}{l}{\textbf{TRENNBARE VERBEN}}\\
& \textbf{---} & \\
\end{tabular}
\vspace{-5pt}
\end{verbentry}

% bewerben (to advertise / promote)
\begin{verbentry}{
  style = {acc},
  toc = {bewerben (to advertise / promote)},
  name = {bewerben},
  english = {to advertise / promote},
  type = {irregular, + Akkusativ},
  auxiliary = {haben}
}


\vspace{5pt}
\hrule
\vspace{5pt}

\begin{tabular}{rl}
\multicolumn{2}{l}{\textbf{PRÄSENS}}\\
ich & \textbf{bewerbe}\\
du & \textbf{bewirbst}\\
er/sie/es & \textbf{bewirbt}\\
wir & \textbf{bewerben}\\
ihr & \textbf{bewerbt}\\
sie/Sie & \textbf{bewerben}\\
\end{tabular}
\hfill
\begin{tabular}{rl}
\multicolumn{2}{l}{\textbf{PRÄTERITUM}}\\
ich & \textbf{bewarb}\\
du & \textbf{bewarbst}\\
er/sie/es & \textbf{bewarb}\\
wir & \textbf{bewarben}\\
ihr & \textbf{bewarbt}\\
sie/Sie & \textbf{bewarben}\\
\end{tabular}
\hfill
\begin{tabular}{rl}
\multicolumn{2}{l}{\textbf{IMPERATIV}}\\
& \textbf{bewirb (du)!}\\
& \textbf{bewerben wir!}\\
& \textbf{bewerbt (ihr)!}\\
& \textbf{bewerben Sie!}\\
\end{tabular}

\vspace{10pt}

\begin{tabular}{rl}
\multicolumn{2}{l}{\textbf{PERFEKT}}\\
& haben + \textbf{beworben}\\
\end{tabular}
\hfill
\begin{tabular}{rl}
\multicolumn{2}{l}{\textbf{FUTURE I}}\\
& werden + \textbf{bewerben}\\
\end{tabular}
\hfill
\begin{tabular}{rl}
\multicolumn{2}{l}{\textbf{PLUSQUAMPERFEKT}}\\
& hatten + \textbf{beworben}\\
\end{tabular}

\vspace{10pt}

\begin{tabular}{lll}
\multicolumn{3}{l}{\textbf{MIT PRÄPOSITIONEN}}\\
& \textbf{---} & \\
\end{tabular}
\hfill
\begin{tabular}{lll}
\multicolumn{3}{l}{\textbf{TRENNBARE VERBEN}}\\
& \textbf{---} & \\
\end{tabular}

\vspace{-5pt}
\end{verbentry}

% bezahlen (to pay)
\begin{verbentry}{
  style = {acc},
  toc = {bezahlen (to pay)},
  name = {bezahlen},
  english = {to pay},
  type = {regular, + Akkusativ},
  auxiliary = {haben}
}
 
    
     \vspace{5pt}

    \hrule
    
    \vspace{5pt}

  \begin{tabular}{rl}
     \multicolumn{2}{l}{\textbf{PRÄSENS} }\\
    ich & \textbf{bezahle}\\
    du & \textbf{bezahlst}\\
    er/sie/es & \textbf{bezahlt}\\
    wir & \textbf{bezahlen}\\
    ihr & \textbf{bezahlt}\\
    sie/Sie & \textbf{bezahlen}\\
  \end{tabular}
  \hfill
     \begin{tabular}{rl}
    \multicolumn{2}{l}{\textbf{PRÄTERITUM} }\\
    ich & \textbf{bezahlte}\\
    du & \textbf{bezahltest}\\
    er/sie/es & \textbf{bezahlte}\\
    wir & \textbf{bezahlten}\\
    ihr & \textbf{bezahltet}\\
    sie/Sie & \textbf{bezahlten}\\
  \end{tabular}
  \hfill
  \begin{tabular}{rl}
     \multicolumn{2}{l}{\textbf{IMPERATIV} }\\
   & \textbf{bezahl(e) (du)!}\\
   & \textbf{bezahlen wir!}\\
   & \textbf{bezahlt (ihr)!}\\
   & \textbf{bezahlen Sie!}\\
   & \vspace{10pt}
  \end{tabular}

    \vspace{10pt}

 \begin{tabular}{rl}
     \multicolumn{2}{l}{\textbf{PERFEKT}}\\
  &  haben + \textbf{bezahlt}\\
  \end{tabular}
\hfill
   \begin{tabular}{rl}
   
    \multicolumn{2}{l}{\textbf{FUTURE I} }\\
  &  werden + \textbf{bezahlen}\\
 
  \end{tabular}
\hfill
   \begin{tabular}{rl}
      \multicolumn{2}{l}{\textbf{PLUSQUAMPERFEKT}}\\
  &  hatten + \textbf{bezahlt}\\
  \end{tabular}
  
  \vspace{10pt}

\begin{tabular}{lll}
\multicolumn{3}{l}{\textbf{MIT PRÄPOSITIONEN}}\\
& \textbf{---} & \\
\end{tabular}
\hfill
\begin{tabular}{lll}
\multicolumn{3}{l}{\textbf{TRENNBARE VERBEN}}\\
& \textbf{---} & \\
\end{tabular}

\vspace{-5pt}
\end{verbentry}

% beziehen (to obtain / refer to)
\begin{verbentry}{
  style = {acc},
  toc = {beziehen (to obtain / refer to)},
  name = {beziehen},
  english = {to obtain / refer to},
  type = {irregular, + Akkusativ},
  auxiliary = {haben}
}


\vspace{5pt}
\hrule
\vspace{5pt}

\begin{tabular}{rl}
\multicolumn{2}{l}{\textbf{PRÄSENS}}\\
ich & \textbf{beziehe}\\
du & \textbf{beziehst}\\
er/sie/es & \textbf{bezieht}\\
wir & \textbf{beziehen}\\
ihr & \textbf{bezieht}\\
sie/Sie & \textbf{beziehen}\\
\end{tabular}
\hfill
\begin{tabular}{rl}
\multicolumn{2}{l}{\textbf{PRÄTERITUM}}\\
ich & \textbf{bezog}\\
du & \textbf{bezogst}\\
er/sie/es & \textbf{bezog}\\
wir & \textbf{bezogen}\\
ihr & \textbf{bezogt}\\
sie/Sie & \textbf{bezogen}\\
\end{tabular}
\hfill
\begin{tabular}{rl}
\multicolumn{2}{l}{\textbf{IMPERATIV}}\\
& \textbf{bezieh(e) (du)!}\\
& \textbf{beziehen wir!}\\
& \textbf{bezieht (ihr)!}\\
& \textbf{beziehen Sie!}\\
\end{tabular}

\vspace{10pt}

\begin{tabular}{rl}
\multicolumn{2}{l}{\textbf{PERFEKT}}\\
& haben + \textbf{bezogen}\\
\end{tabular}
\hfill
\begin{tabular}{rl}
\multicolumn{2}{l}{\textbf{FUTURE I}}\\
& werden + \textbf{beziehen}\\
\end{tabular}
\hfill
\begin{tabular}{rl}
\multicolumn{2}{l}{\textbf{PLUSQUAMPERFEKT}}\\
& hatten + \textbf{bezogen}\\
\end{tabular}

\vspace{10pt}

\begin{tabular}{lll}
\multicolumn{3}{l}{\textbf{MIT PRÄPOSITIONEN}}\\
& \textbf{beziehen Akk} & (to obtain something)\\
& \textbf{beziehen auf} + Akk & (to refer to)\\
\end{tabular}
\hfill
\begin{tabular}{lll}
\multicolumn{3}{l}{\textbf{TRENNBARE VERBEN}}\\
& \textbf{ein$|$beziehen} & (to involve / include)\\
& \textbf{bevor$|$ziehen} & (to prefer)\\
\end{tabular}

\vspace{-5pt}
\end{verbentry}

% bieten (to offer)
\begin{verbentry}{
  style = {default},
  toc = {bieten (to offer)},
  name = {bieten},
  english = {to offer},
  type = {irregular},
  auxiliary = {haben}
}


\vspace{5pt}
\hrule
\vspace{5pt}

\begin{tabular}{rl}
\multicolumn{2}{l}{\textbf{PRÄSENS}}\\
ich & \textbf{biete}\\
du & \textbf{bietest}\\
er/sie/es & \textbf{bietet}\\
wir & \textbf{bieten}\\
ihr & \textbf{bietet}\\
sie/Sie & \textbf{bieten}\\
\end{tabular}
\hfill
\begin{tabular}{rl}
\multicolumn{2}{l}{\textbf{PRÄTERITUM}}\\
ich & \textbf{bot}\\
du & \textbf{botest}\\
er/sie/es & \textbf{bot}\\
wir & \textbf{boten}\\
ihr & \textbf{botet}\\
sie/Sie & \textbf{boten}\\
\end{tabular}
\hfill
\begin{tabular}{rl}
\multicolumn{2}{l}{\textbf{IMPERATIV}}\\
& \textbf{biete (du)!}\\
& \textbf{bieten wir!}\\
& \textbf{bietet (ihr)!}\\
& \textbf{bieten Sie!}\\
\end{tabular}

\vspace{10pt}

\begin{tabular}{rl}
\multicolumn{2}{l}{\textbf{PERFEKT}}\\
& haben + \textbf{geboten}\\
\end{tabular}
\hfill
\begin{tabular}{rl}
\multicolumn{2}{l}{\textbf{FUTURE I}}\\
& werden + \textbf{bieten}\\
\end{tabular}
\hfill
\begin{tabular}{rl}
\multicolumn{2}{l}{\textbf{PLUSQUAMPERFEKT}}\\
& hatten + \textbf{geboten}\\
\end{tabular}

\vspace{10pt}

\begin{tabular}{lll}
\multicolumn{3}{l}{\textbf{MIT PRÄPOSITIONEN}}\\
& \textbf{bieten} + Dat + Akk & (to offer someone something)\\
& \textbf{bieten für} + Akk & (to bid for)\\
\end{tabular}
\hfill
\begin{tabular}{lll}
\multicolumn{3}{l}{\textbf{TRENNBARE VERBEN}}\\
& \textbf{an$|$bieten} & (to offer)\\
\vspace{10pt}
\end{tabular}

\vspace{-5pt}
\end{verbentry}

% bitten (to ask / request)
\begin{verbentry}{
  style = {default},
  toc = {bitten (to ask / request)},
  name = {bitten},
  english = {to ask / request},
  type = {irregular},
  auxiliary = {haben}
}


\vspace{5pt}
\hrule
\vspace{5pt}

\begin{tabular}{rl}
\multicolumn{2}{l}{\textbf{PRÄSENS}}\\
ich & \textbf{bitte}\\
du & \textbf{bittest}\\
er/sie/es & \textbf{bittet}\\
wir & \textbf{bitten}\\
ihr & \textbf{bittet}\\
sie/Sie & \textbf{bitten}\\
\end{tabular}
\hfill
\begin{tabular}{rl}
\multicolumn{2}{l}{\textbf{PRÄTERITUM}}\\
ich & \textbf{bat}\\
du & \textbf{batst}\\
er/sie/es & \textbf{bat}\\
wir & \textbf{baten}\\
ihr & \textbf{batet}\\
sie/Sie & \textbf{baten}\\
\end{tabular}
\hfill
\begin{tabular}{rl}
\multicolumn{2}{l}{\textbf{IMPERATIV}}\\
& \textbf{bitte (du)!}\\
& \textbf{bitten wir!}\\
& \textbf{bittet (ihr)!}\\
& \textbf{bitten Sie!}\\
\end{tabular}

\vspace{10pt}

\begin{tabular}{rl}
\multicolumn{2}{l}{\textbf{PERFEKT}}\\
& haben + \textbf{gebeten}\\
\end{tabular}
\hfill
\begin{tabular}{rl}
\multicolumn{2}{l}{\textbf{FUTURE I}}\\
& werden + \textbf{bitten}\\
\end{tabular}
\hfill
\begin{tabular}{rl}
\multicolumn{2}{l}{\textbf{PLUSQUAMPERFEKT}}\\
& hatten + \textbf{gebeten}\\
\end{tabular}

\vspace{10pt}

\begin{tabular}{lll}
\multicolumn{3}{l}{\textbf{MIT PRÄPOSITIONEN}}\\
& \textbf{bitten um} + Akk & (to ask for / request)\\
\end{tabular}
\hfill
\begin{tabular}{lll}
\multicolumn{3}{l}{\textbf{TRENNBARE VERBEN}}\\
& \textbf{---} & \\
\end{tabular}
\vspace{-5pt}
\end{verbentry}

% bleiben (to stay)
\begin{verbentry}{
  style = {default},
  toc = {bleiben (to stay)},
  name = {bleiben},
  english = {to stay},
  type = {irregular},
  auxiliary = {sein}
}
 
    
     \vspace{5pt}

    \hrule
    
    \vspace{5pt}

  \begin{tabular}{rl}
     \multicolumn{2}{l}{\textbf{PRÄSENS} }\\
    ich & \textbf{bleibe}\\
    du & \textbf{bleibst}\\
    er/sie/es & \textbf{bleibt}\\
    wir & \textbf{bleiben}\\
    ihr & \textbf{bleibt}\\
    sie/Sie & \textbf{bleiben}\\
  \end{tabular}
  \hfill
     \begin{tabular}{rl}
    \multicolumn{2}{l}{\textbf{PRÄTERITUM} }\\
    ich & \textbf{blieb}\\
    du & \textbf{bliebst}\\
    er/sie/es & \textbf{blieb}\\
    wir & \textbf{blieben}\\
    ihr & \textbf{bliebt}\\
    sie/Sie & \textbf{blieben}\\
  \end{tabular}
  \hfill
  \begin{tabular}{rl}
     \multicolumn{2}{l}{\textbf{IMPERATIV} }\\
   & \textbf{bleib(e) (du)!}\\
   & \textbf{bleiben wir!}\\
   & \textbf{bleibt (ihr)!}\\
   & \textbf{bleiben Sie!}\\
   & \vspace{10pt}
  \end{tabular}

    \vspace{10pt}

 \begin{tabular}{rl}
     \multicolumn{2}{l}{\textbf{PERFEKT}}\\
  &  sein + \textbf{geblieben}\\
  \end{tabular}
\hfill
   \begin{tabular}{rl}
   
    \multicolumn{2}{l}{\textbf{FUTURE I} }\\
  &  werden + \textbf{bleiben}\\
 
  \end{tabular}
\hfill
   \begin{tabular}{rl}
      \multicolumn{2}{l}{\textbf{PLUSQUAMPERFEKT}}\\
  &  waren + \textbf{geblieben}\\
  \end{tabular}
  
  \vspace{10pt}

\begin{tabular}{lll}
\multicolumn{3}{l}{\textbf{MIT PRÄPOSITIONEN}}\\
& \textbf{---} & \\
\end{tabular}
\hfill
\begin{tabular}{lll}
\multicolumn{3}{l}{\textbf{TRENNBARE VERBEN}}\\
& \textbf{---} & \\
\end{tabular}

  \vspace{-5pt}
\end{verbentry}

% blühen (to bloom / flourish)
\begin{verbentry}{
  style = {default},
  toc = {blühen (to bloom / flourish)},
  name = {blühen},
  english = {to bloom / flourish},
  type = {regular},
  auxiliary = {haben}
}


\vspace{5pt}
\hrule
\vspace{5pt}

\begin{tabular}{rl}
\multicolumn{2}{l}{\textbf{PRÄSENS}}\\
ich & \textbf{blühe}\\
du & \textbf{blühst}\\
er/sie/es & \textbf{blüht}\\
wir & \textbf{blühen}\\
ihr & \textbf{blüht}\\
sie/Sie & \textbf{blühen}\\
\end{tabular}
\hfill
\begin{tabular}{rl}
\multicolumn{2}{l}{\textbf{PRÄTERITUM}}\\
ich & \textbf{blühte}\\
du & \textbf{blühtest}\\
er/sie/es & \textbf{blühte}\\
wir & \textbf{blühten}\\
ihr & \textbf{blühtet}\\
sie/Sie & \textbf{blühten}\\
\end{tabular}
\hfill
\begin{tabular}{rl}
\multicolumn{2}{l}{\textbf{IMPERATIV}}\\
& \textbf{blühe (du)!}\\
& \textbf{blühen wir!}\\
& \textbf{blüht (ihr)!}\\
& \textbf{blühen Sie!}\\
\end{tabular}

\vspace{10pt}

\begin{tabular}{rl}
\multicolumn{2}{l}{\textbf{PERFEKT}}\\
& haben + \textbf{geblüht}\\
\end{tabular}
\hfill
\begin{tabular}{rl}
\multicolumn{2}{l}{\textbf{FUTURE I}}\\
& werden + \textbf{blühen}\\
\end{tabular}
\hfill
\begin{tabular}{rl}
\multicolumn{2}{l}{\textbf{PLUSQUAMPERFEKT}}\\
& hatten + \textbf{geblüht}\\
\end{tabular}

\vspace{10pt}

\begin{tabular}{lll}
\multicolumn{3}{l}{\textbf{MIT PRÄPOSITIONEN}}\\
& \textbf{---} & \\
\end{tabular}
\hfill
\begin{tabular}{lll}
\multicolumn{3}{l}{\textbf{TRENNBARE VERBEN}}\\
& \textbf{auf$|$blühen} & (to blossom / flourish)\\
\end{tabular}
\vspace{-5pt}
\end{verbentry}

% bohren (to drill)
\begin{verbentry}{
  style = {default},
  toc = {bohren (to drill)},
  name = {bohren},
  english = {to drill},
  type = {regular},
  auxiliary = {haben}
}


\vspace{5pt}
\hrule
\vspace{5pt}

\begin{tabular}{rl}
\multicolumn{2}{l}{\textbf{PRÄSENS}}\\
ich & \textbf{bohre}\\
du & \textbf{bohrst}\\
er/sie/es & \textbf{bohrt}\\
wir & \textbf{bohren}\\
ihr & \textbf{bohrt}\\
sie/Sie & \textbf{bohren}\\
\end{tabular}
\hfill
\begin{tabular}{rl}
\multicolumn{2}{l}{\textbf{PRÄTERITUM}}\\
ich & \textbf{bohrte}\\
du & \textbf{bohrtest}\\
er/sie/es & \textbf{bohrte}\\
wir & \textbf{bohrten}\\
ihr & \textbf{bohrtet}\\
sie/Sie & \textbf{bohrten}\\
\end{tabular}
\hfill
\begin{tabular}{rl}
\multicolumn{2}{l}{\textbf{IMPERATIV}}\\
& \textbf{bohre (du)!}\\
& \textbf{bohren wir!}\\
& \textbf{bohrt (ihr)!}\\
& \textbf{bohren Sie!}\\
\end{tabular}

\vspace{10pt}

\begin{tabular}{rl}
\multicolumn{2}{l}{\textbf{PERFEKT}}\\
& haben + \textbf{gebohrt}\\
\end{tabular}
\hfill
\begin{tabular}{rl}
\multicolumn{2}{l}{\textbf{FUTURE I}}\\
& werden + \textbf{bohren}\\
\end{tabular}
\hfill
\begin{tabular}{rl}
\multicolumn{2}{l}{\textbf{PLUSQUAMPERFEKT}}\\
& hatten + \textbf{gebohrt}\\
\end{tabular}

\vspace{10pt}

\begin{tabular}{lll}
\multicolumn{3}{l}{\textbf{MIT PRÄPOSITIONEN}}\\
& \textbf{bohren in} + Dat & (to drill in / into)\\
& \textbf{bohren durch} + Akk & (to drill through)\\
\end{tabular}
\hfill
\begin{tabular}{lll}
\multicolumn{3}{l}{\textbf{TRENNBARE VERBEN}}\\
& \textbf{auf$|$bohren} & (to drill on / onto surface)\\
\end{tabular}
\vspace{-5pt}
\end{verbentry}

% brauchen (to need)
\begin{verbentry}{
  style = {acc},
  toc = {brauchen (to need)},
  name = {brauchen},
  english = {to need},
  type = {regular, + Akkusativ},
  auxiliary = {haben}
}
 
    
     \vspace{5pt}

    \hrule
    
    \vspace{5pt}

  \begin{tabular}{rl}
     \multicolumn{2}{l}{\textbf{PRÄSENS} }\\
    ich & \textbf{brauche}\\
    du & \textbf{brauchst}\\
    er/sie/es & \textbf{braucht}\\
    wir & \textbf{brauchen}\\
    ihr & \textbf{braucht}\\
    sie/Sie & \textbf{brauchen}\\
  \end{tabular}
  \hfill
     \begin{tabular}{rl}
    \multicolumn{2}{l}{\textbf{PRÄTERITUM} }\\
    ich & \textbf{brauchte}\\
    du & \textbf{brauchtest}\\
    er/sie/es & \textbf{brauchte}\\
    wir & \textbf{brauchten}\\
    ihr & \textbf{brauchtet}\\
    sie/Sie & \textbf{brauchten}\\
  \end{tabular}
  \hfill
  \begin{tabular}{rl}
     \multicolumn{2}{l}{\textbf{IMPERATIV} }\\
   & \textbf{brauch(e) (du)!}\\
   & \textbf{brauchen wir!}\\
   & \textbf{braucht (ihr)!}\\
   & \textbf{brauchen Sie!}\\
   & \vspace{10pt}
  \end{tabular}

    \vspace{10pt}

 \begin{tabular}{rl}
     \multicolumn{2}{l}{\textbf{PERFEKT}}\\
   & haben + \textbf{gebraucht}\\
  \end{tabular}
\hfill
   \begin{tabular}{rl}
   
    \multicolumn{2}{l}{\textbf{FUTURE I} }\\
  &  werden + \textbf{brauchen}\\
 
  \end{tabular}
\hfill
   \begin{tabular}{rl}
      \multicolumn{2}{l}{\textbf{PLUSQUAMPERFEKT}}\\
  &  hatten + \textbf{gebraucht}\\
  \end{tabular}
  
  \vspace{10pt}

\begin{tabular}{lll}
\multicolumn{3}{l}{\textbf{MIT PRÄPOSITIONEN}}\\
& \textbf{---} & \\
\end{tabular}
\hfill
\begin{tabular}{lll}
\multicolumn{3}{l}{\textbf{TRENNBARE VERBEN}}\\
& \textbf{---} & \\
\end{tabular}

  \vspace{-5pt}
\end{verbentry}

% bringen (to bring)
\begin{verbentry}{
  style = {acc},
  toc = {bringen (to bring)},
  name = {bringen},
  english = {to bring},
  type = {irregular, + Akkusativ},
  auxiliary = {haben}
}
 
    
     \vspace{5pt}

    \hrule
    
    \vspace{5pt}

  \begin{tabular}{rl}
     \multicolumn{2}{l}{\textbf{PRÄSENS} }\\
    ich & \textbf{bringe}\\
    du & \textbf{bringst}\\
    er/sie/es & \textbf{bringt}\\
    wir & \textbf{bringen}\\
    ihr & \textbf{bringt}\\
    sie/Sie & \textbf{bringen}\\
  \end{tabular}
  \hfill
     \begin{tabular}{rl}
    \multicolumn{2}{l}{\textbf{PRÄTERITUM} }\\
    ich & \textbf{brachte}\\
    du & \textbf{brachtest}\\
    er/sie/es & \textbf{brachte}\\
    wir & \textbf{brachten}\\
    ihr & \textbf{brachtet}\\
    sie/Sie & \textbf{brachten}\\
  \end{tabular}
  \hfill
  \begin{tabular}{rl}
     \multicolumn{2}{l}{\textbf{IMPERATIV} }\\
   & \textbf{bring(e) (du)!}\\
   & \textbf{bringen wir!}\\
   & \textbf{bringt (ihr)!}\\
   & \textbf{bringen Sie!}\\
   & \vspace{10pt}
  \end{tabular}

    \vspace{10pt}

 \begin{tabular}{rl}
     \multicolumn{2}{l}{\textbf{PERFEKT}}\\
 &   haben + \textbf{gebracht}\\
  \end{tabular}
\hfill
   \begin{tabular}{rl}
   
    \multicolumn{2}{l}{\textbf{FUTURE I} }\\
  &  werden + \textbf{bringen}\\
 
  \end{tabular}
\hfill
   \begin{tabular}{rl}
      \multicolumn{2}{l}{\textbf{PLUSQUAMPERFEKT}}\\
   & hatten + \textbf{gebracht}\\
  \end{tabular}

  
  \vspace{10pt}

   \begin{tabular}{lll}
      \multicolumn{3}{l}{\textbf{MIT PRÄPOSITIONEN}}\\
& \textbf{bringen zu} + Dat & (to bring to someone / place)\\
&\textbf{bringen auf} + Akk & (to raise / summon)\\
&\textbf{bringen für} + Akk & (to bring for)\\
&\textbf{bringen von} + Dat & (to bring from)\\
&\textbf{bringen in} + Akk & (to move into)\\
\vspace{70pt}
  \end{tabular}
  \hfill
   \begin{tabular}{lll}
      \multicolumn{3}{l}{\textbf{TRENNBARE VERBEN}}\\
 & \textbf{um$|$bringen} & (to kill)\\
 & \textbf{mit$|$bringen} & (to bring along)\\
  & \textbf{nach$|$bringen} & (to bring later)\\
  & \textbf{hin$|$bringen} & (to bring there)\\
  & \textbf{her$|$bringen} & (to bring here)\\
     & \textbf{auf$|$bringen} & (to raise money / summon)\\
     & \textbf{ein$|$bringen} & (to yield / introduce)\\
     & \textbf{bei$|$bringen} & (to teach / instill)\\
     & \textbf{unter$|$bringen} & (to accomodate)\\
     & \textbf{weg$|$bringen} & (to take away)\\
      & \textbf{zurück$|$bringen} & (to bring back)\\

      & \textbf{vorbei$|$bringen} & (to bring by / drop off)\\
  \end{tabular}


  \vspace{-5pt}
\end{verbentry}

% buchen (to book / reserve)
\begin{verbentry}{
  style = {acc},
  toc = {buchen (to book / reserve)},
  name = {buchen},
  english = {to book / reserve},
  type = {regular, + Akkusativ},
  auxiliary = {haben}
}


\vspace{5pt}
\hrule
\vspace{5pt}

\begin{tabular}{rl}
\multicolumn{2}{l}{\textbf{PRÄSENS}}\\
ich & \textbf{buche}\\
du & \textbf{buchst}\\
er/sie/es & \textbf{bucht}\\
wir & \textbf{buchen}\\
ihr & \textbf{bucht}\\
sie/Sie & \textbf{buchen}\\
\end{tabular}
\hfill
\begin{tabular}{rl}
\multicolumn{2}{l}{\textbf{PRÄTERITUM}}\\
ich & \textbf{buchte}\\
du & \textbf{buchtest}\\
er/sie/es & \textbf{buchte}\\
wir & \textbf{buchten}\\
ihr & \textbf{buchtet}\\
sie/Sie & \textbf{buchten}\\
\end{tabular}
\hfill
\begin{tabular}{rl}
\multicolumn{2}{l}{\textbf{IMPERATIV}}\\
& \textbf{buche (du)!}\\
& \textbf{buchen wir!}\\
& \textbf{bucht (ihr)!}\\
& \textbf{buchen Sie!}\\
\end{tabular}

\vspace{10pt}

\begin{tabular}{rl}
\multicolumn{2}{l}{\textbf{PERFEKT}}\\
& haben + \textbf{gebucht}\\
\end{tabular}
\hfill
\begin{tabular}{rl}
\multicolumn{2}{l}{\textbf{FUTURE I}}\\
& werden + \textbf{buchen}\\
\end{tabular}
\hfill
\begin{tabular}{rl}
\multicolumn{2}{l}{\textbf{PLUSQUAMPERFEKT}}\\
& hatten + \textbf{gebucht}\\
\end{tabular}

\vspace{10pt}

\begin{tabular}{lll}
\multicolumn{3}{l}{\textbf{MIT PRÄPOSITIONEN}}\\
& \textbf{buchen für} + Akk & (to book for)\\
& \textbf{buchen bei} + Dat & (to book at / with)\\
\end{tabular}
\hfill
\begin{tabular}{lll}
\multicolumn{3}{l}{\textbf{TRENNBARE VERBEN}}\\
& \textbf{---} & \\
\end{tabular}
\vspace{-5pt}
\end{verbentry}
